\color{unidarkgreen}{\subsection{Дифференциальная защита шин}}
\color{black}

\begin{enumerate}[label=\arabic{section}.\arabic{subsection}.\arabic{enumi}, labelsep=4pt, leftmargin=0pt, itemindent=57pt, itemsep=0pt, parsep=5pt] % если что можно крутить itemsep

%%%%%%%%%%%%%%%%%%%%%%%%%%%%%%%%%%%%%%%%%%%%%%%%%% ОБШИЕ СВЕДЕНИЯ %%%%%%%%%%%%%%%%%%%%%%%%%%%%%%%%%%%%%%%%%%%%%%%%%%
\item Общие сведения

\begin{enumerate}[label=\arabic{section}.\arabic{subsection}.\arabic{enumi}.\arabic*, labelsep=4pt, leftmargin=0em, itemindent=65pt, parsep=0pt]

\item
ДЗШ является быстродействующей защитой с абсолютной селективностью и выполняет функцию основной токовой защиты шин от всех видов КЗ.

\item
Данное устройство ДЗШ имеет пофазное исполнение, т.е. каждая фаза обрабатывается одним терминалом независимо от двух других.
Для селективной защиты одной фазы устройство содержит три дифференциальных органа: пусковой орган (ПО) для защиты обеих СШ, избирательные органы для защиты первой (ИО1) и второй (ИО2) СШ.
Ввод и вывод данной функции для всех дифференциальных органов осуществляется либо посредством одного программного ключа \textbf{«ДЗШ»}, либо с использованием внешнего сигнала \textit{«Вывод ДЗШ»}.

\item
При рабочей фиксации присоединений сигнал на отключение поврежденной СШ действует только при срабатывании пускового органа и соответствующего избирательного органа.
В режиме нарушенной фиксации избирательные органы выводятся из работы и отключение обеих СШ производится только от пускового органа.

\item
Пусковой орган подключен на токи всех присоединений обеих СШ, за исключением токов ШСВ.
Избирательные органы подключены на токи всех присоединений соответствующей СШ, включая токи ШСВ.
Логика фиксации присоединений приведена в разделе «Логика управления топологией».

\item
В устройстве предусмотрена настройка полярности ТТ присоединений.
По умолчанию измерение токов выполняется в направлении к защищаемому объекту, то есть все ТТ считаются включёнными встречно друг к другу.
Если для какого-либо присоединения фактическая полярность не совпадает с принятой по умолчанию, пользователь может скорректировать её двумя способами: либо изменением подключения вторичных токовых цепей, либо установкой программного ключа \textbf{«Инверсия Bn»} в положение \textbf{«Введено»}.

\small
\begin{longtable}{|>{\centering\arraybackslash}m{6.4cm}|>{\centering\arraybackslash}m{4cm}|>{\centering\arraybackslash}m{2.8cm}|>{\centering\arraybackslash}m{2.8cm}|}
% Первая страница: полный заголовок
\caption{Перечень параметров ТТ и ТН\hfill\vspace{-0.5\baselineskip}}\label{desc:tbl1}\\
\hline
\rowcolor{gray!30}
Наименование уставки & Диапазон & Шаг & Значение по умолчанию \\
\hline
\endfirsthead

% Промежуточные страницы: «Продолжение таблицы»
\caption*{\hspace{3pt}\emph{Продолжение таблицы \ref{desc:tbl1}}\hfill\ }\vspace{-0.5\baselineskip} \\
\hline
\rowcolor{gray!30}
Наименование уставки & Диапазон & Шаг & Значение по умолчанию \\
\hline
\endhead

\centering\arraybackslash Номинальный первичный ток ТТ В1, А  & \centering\arraybackslash 1…12000 & \centering\arraybackslash 1 & \centering\arraybackslash 1000 \\
\hline
\centering\arraybackslash Номинальный вторичный ток ТТ В1, А  & \centering\arraybackslash 1…12000 & \centering\arraybackslash 1 & \centering\arraybackslash 1000 \\
\hline
\centering\arraybackslash Инверсия В1  & \centering Введено \\ Выведено & \centering --- & \centering\arraybackslash Выведено \\
\hline
\centering\arraybackslash …  & \centering\arraybackslash … & \centering\arraybackslash … & \centering\arraybackslash … \\
\hline
\centering\arraybackslash Номинальный первичный ток ТТ В12, А  & \centering\arraybackslash 1…12000 & \centering\arraybackslash 1 & \centering\arraybackslash 1000 \\
\hline
\centering\arraybackslash Номинальный вторичный ток ТТ В12, А  & \centering\arraybackslash 1…12000 & \centering\arraybackslash 1 & \centering\arraybackslash 1000 \\
\hline
\centering\arraybackslash Инверсия В12  & \centering Введено \\ Выведено & \centering --- & \centering\arraybackslash Выведено \\
\hline
\centering\arraybackslash Номинальное первичное линейное напряжение ТН 1С, кВ  & \centering\arraybackslash 0,22…1150,00 & \centering\arraybackslash 0,01 & \centering\arraybackslash 110 \\
\hline
\centering\arraybackslash Номинальное вторичное линейное напряжение ТН 1С, кВ  & \centering\arraybackslash 100,00…400,00 & \centering\arraybackslash 0,01 & \centering\arraybackslash 100 \\
\hline
\centering\arraybackslash Номинальное первичное линейное напряжение ТН 2С, кВ  & \centering\arraybackslash 0,22…1150,00 & \centering\arraybackslash 0,01 & \centering\arraybackslash 110 \\
\hline
\centering\arraybackslash Номинальное вторичное линейное напряжение ТН 2С, кВ  & \centering\arraybackslash 100,00…400,00 & \centering\arraybackslash 0,01 & \centering\arraybackslash 100 \\
\hline

\end{longtable}
\normalsize

\end{enumerate}


%%%%%%%%%%%%%%%%%%%%%%%%%%%%%%%%%%%%%%%%%%%%%%%%%% Цифровое выравнивание токов %%%%%%%%%%%%%%%%%%%%%%%%%%%%%%%%%%%%%%%%%%%%%%%%%%
\item Цифровое выравнивание токов

\begin{enumerate}[label=\arabic{section}.\arabic{subsection}.\arabic{enumi}.\arabic*, labelsep=4pt, leftmargin=0em, itemindent=65pt, parsep=0pt]

\item
В устройстве реализовано цифровое выравнивание токов присоединений, что позволяет подключать защиту к цепям ТТ с различными коэффициентами трансформации без использования промежуточных преобразователей.

\item
Токи присоединений приводятся к общему базисному току по выражению:
\begin{align}
i_n^{'}=\frac{i_n}{I_\textit{баз}}
\end{align}
\begin{eqexpl} [2mm]
    \item{$ i_n^{'} $} мгновенное значение выровненного тока n-го присоединения, о.е.;
	\item{$ i_n $} мгновенное значение первичного тока n-го присоединения, А;
	\item{$ I_\textit{баз} $} базисный ток, А.
\end{eqexpl}

\item
Величина базисного тока задается уставкой «Базисный ток» В качестве базисного тока выбирается наибольший номинальный первичный ток ТТ из всех присоединений:
\begin{align}
I_\textit{баз}=max(I_\textit{перв.ном,1}...I_\textit{перв.ном,n})
\end{align}

\end{enumerate}


%%%%%%%%%%%%%%%%%%%%%%%%%%%%%%%%%%%%%%%%%%%%%%%%%% Измерительные органы ДЗШ %%%%%%%%%%%%%%%%%%%%%%%%%%%%%%%%%%%%%%%%%%%%%%%%%%
\item Измерительные органы ДЗШ

\begin{enumerate}[label=\arabic{section}.\arabic{subsection}.\arabic{enumi}.\arabic*, labelsep=4pt, leftmargin=0em, itemindent=65pt, parsep=0pt]

\item
Каждый дифференциальный орган имеет идентичную реализацию и состоит из следующих измерительных органов:
\begin{itemize} 
\item дифференциальная токовая отсечка (ДТО);
\item дифференциальный орган с торможением (ДОТХ);
\item быстродействующий дифференциальный орган (БДО);
\item медленнодействующий дифференциальный орган (МДО);
\item блокировка при внешних коротких замыканиях (БВКЗ);
\item чувствительный токовый орган (ЧТО);
\item орган блокировки по второй гармонике дифференциального тока;
\item орган дифференциального тока для КЦТ.
\end{itemize}

\medskip

\begin{figure}[H]
\centering
\includegraphics[width=0.8\textwidth,height=0.8\textheight,keepaspectratio]{\fbpath/661.0631 ДЗШ/img/IO_DZSH.pdf}
\captionof{figure}{Функциональный блок ДЗШ ИО1}\label{pict:IO_DZSH}
\end{figure}

\medskip

\end{enumerate}


%%%%%%%%%%%%%%%%%%%%%%%%%%%%%%%%%%%%%%%%%%%%%%%%%% Дифференциальная токовая отсечка %%%%%%%%%%%%%%%%%%%%%%%%%%%%%%%%%%%%%%%%%%%%%%%%%%
\item Дифференциальная токовая отсечка

\begin{enumerate}[label=\arabic{section}.\arabic{subsection}.\arabic{enumi}.\arabic*, labelsep=4pt, leftmargin=0em, itemindent=65pt, parsep=0pt]

\item
ДТО предназначена для быстрого отключения защищаемого объекта при внутренних повреждениях, сопровождающихся большими токами.
Ввод в работу данной функции для всех дифференциальных органов осуществляется с помощью одного программного ключа \textbf{«ДТО»}.

\item
Данный ИО выполнен без торможения и может быть заблокирован только при обнаружении неисправности в ЦТ.
В качестве рабочей величины использует действующее значение первой гармоники дифференциального тока, которое определяется как модуль геометрической суммы первых гармоник токов присоединений, входящих в зону действия защиты:
\begin{align}
I_\textit{Д}=\mid \sum_{n=1}^N \vec{I}_{n}^{\,'} \mid
\end{align}
\begin{eqexpl} [2mm]
    \item{$ I_\textit{Д} $} действующее значение первой гармоники дифференциального тока, о.е.;
	\item{$ N $} количество присоединений, зафиксированных за соответствующим дифференциальным органом;
	\item{$ \vec{I}_{n}^{\,'} $} вектор первой гармоники выровненного тока n-го присоединения, о.е..
\end{eqexpl}

\item
Ток срабатывания ДТО для каждого дифференциального органа задается уставкой \textbf{«Iдто ПО»} (\textbf{«Iдто ИО1», «Iдто ИО2»}).


\end{enumerate}


%%%%%%%%%%%%%%%%%%%%%%%%%%%%%%%%%%%%%%%%%%%%%%%%%% Дифференциальный орган с торможением %%%%%%%%%%%%%%%%%%%%%%%%%%%%%%%%%%%%%%%%%%%%%%%%%%
\item Дифференциальный орган с торможением

\begin{enumerate}[label=\arabic{section}.\arabic{subsection}.\arabic{enumi}.\arabic*, labelsep=4pt, leftmargin=0em, itemindent=65pt, parsep=0pt]

\item
ДОТХ является основным измерительным органом и предназначен для отключения защищаемого объекта при всех видах внутренних повреждений.

\item
В качестве рабочих величин данный ИО использует действующие значения дифференциального и тормозного токов первой гармоники.
Действующее значение тормозного тока первой гармоники рассчитывается как полусумма модулей токов присоединений, входящих в зону действия защиты:
\begin{align}
I_\textit{Т}=0,5 \cdot \sum_{n=1}^N \mid \vec{I}_{n}^{\,'} \mid
\end{align}

\item
Тормозная характеристика дифференциальной защиты имеет два участка.
Горизонтальный участок не зависит от тормозного тока и необходим для отстройки от небольших токов небаланса при протекании малых сквозных токов, которые не вызывают насыщение ТТ.
Уровень срабатывания данного участка для каждого дифференциального органа задается отдельной уставкой \textbf{«Iд нач. ПО»} (\textbf{«Iд нач. ИО1»}, \textbf{«Iд нач. ИО2»}).
Наклонный участок обеспечивает отстройку от токов небаланса при насыщении ТТ.
Наклон характеристики определяется коэффициентом торможения, а точка начала наклона – значением начального тормозного тока.
Для каждого дифференциального органа эти параметры задаются отдельными уставками: \textbf{«Kт ПО»} и \textbf{«Iт нач. ПО»} (\textbf{«Kт ИО1»} и \textbf{«Iт нач. ИО1»}, \textbf{«Kт ИО2»} и \textbf{«Iт нач. ИО2»}).

\medskip

\begin{figure}[H]
\centering
\includegraphics[width=0.55\textwidth,height=0.55\textheight,keepaspectratio]{\fbpath/661.0631 ДЗШ/img/DZSH_torm_char.pdf}
\captionof{figure}{Тормозная характеристика дифференциальной защиты}\label{pic:DZSH_torm_char}
\end{figure}

\medskip

\end{enumerate}


%%%%%%%%%%%%%%%%%%%%%%%%%%%%%%%%%%%%%%%%%%%%%%%%%% Быстродействующий орган %%%%%%%%%%%%%%%%%%%%%%%%%%%%%%%%%%%%%%%%%%%%%%%%%%
\item Быстродействующий орган

\begin{enumerate}[label=\arabic{section}.\arabic{subsection}.\arabic{enumi}.\arabic*, labelsep=4pt, leftmargin=0em, itemindent=65pt, parsep=0pt]

\item
БДО предназначен для быстрого выявления внутренних повреждений до наступления насыщения ТТ.
Ввод в работу данной функции для всех дифференциальных органов осуществляется с помощью одного программного ключа \textbf{«БДО»}.

\item
Принцип действия органа основан на различии скорости нарастания дифференциального и тормозного токов при внутренних и внешних повреждениях.
При внутреннем КЗ обеспечивается равный рост дифференциального и тормозного тока в промежутке правильной трансформации ТТ.
При внешнем КЗ нарастание тормозного тока происходит раньше, тогда как рост дифференциального тока обусловлен насыщением ТТ.

\item
В алгоритме БДО реализована логика выявления одновременного нарастания фронтов действующих значений дифференциального и тормозного токов.
Сигнал срабатывания БДО формируется при наличии превышении дифференциального тока уставки срабатывания ДОТХ и нарастанием фронтов дифференциального и тормозного токов в пределах заданного времени.

\end{enumerate}


%%%%%%%%%%%%%%%%%%%%%%%%%%%%%%%%%%%%%%%%%%%%%%%%%% Медленнодействующий орган %%%%%%%%%%%%%%%%%%%%%%%%%%%%%%%%%%%%%%%%%%%%%%%%%%
\item Медленнодействующий орган

\begin{enumerate}[label=\arabic{section}.\arabic{subsection}.\arabic{enumi}.\arabic*, labelsep=4pt, leftmargin=0em, itemindent=65pt, parsep=0pt]

\item
МДО предназначен для выявления внутренних повреждений, а также для распознавания перехода КЗ из внешнего во внутреннее, в том числе в условиях насыщения ТТ.
Ввод в работу МДО осуществляется программным ключом \textbf{«МДО»}.

\item
Принцип действия органа основан на контроле формы абсолютного (выпрямленного) мгновенного дифференциального тока.
При внутреннем КЗ мгновенный дифференциальный ток изменяется два раза за период промышленной частоты, в то время как при внешнем КЗ с насыщением ТТ – один раз за период.
МДО обеспечивает требование по быстродействию при переходе внешнего КЗ во внутреннее.

\item
Выпрямленный мгновенный дифференциальный ток рассчитывается как модуль суммы мгновенных значений токов присоединений, входящих в зону действия защиты:
\begin{align}
i_\textit{д}=\mid \sum_{n=1}^N i_{n} \mid,
\end{align}
\begin{eqexpl} [0mm]
    \item{$ i_\textit{д} $} мгновенное значение выпрямленного дифференциального тока, о.е..
\end{eqexpl}

\end{enumerate}


%%%%%%%%%%%%%%%%%%%%%%%%%%%%%%%%%%%%%%%%%%%%%%%%%% Орган блокировки при внешнем КЗ %%%%%%%%%%%%%%%%%%%%%%%%%%%%%%%%%%%%%%%%%%%%%%%%%%
\item Орган блокировки при внешнем КЗ

\begin{enumerate}[label=\arabic{section}.\arabic{subsection}.\arabic{enumi}.\arabic*, labelsep=4pt, leftmargin=0em, itemindent=65pt, parsep=0pt]

\item
БВКЗ предназначена для выявления внешнего КЗ и обеспечения времени достоверного измерения тока на уровне 5 мс.
Ввод в работу данной функции для всех дифференциальных органов осуществляется с помощью одного программного ключа \textbf{«БВКЗ»}.

\item
Принцип действия органа основан на сравнении направлений токов присоединений.
Для идентификации внешнего повреждения формируются два вектора: вектор тока повреждённого присоединения ($ \vec{I}_1 $), вектор тока, рассчитываемого как разность между дифференциальным током и током повреждения ($ \vec{I}_2 $), т.е. вектор суммы токов всех остальных присоединений.

\item
В качестве вектора тока поврежденного присоединения выступает комплексное действующее значение основной гармоники тока того присоединения, в котором зафиксировано максимальное значение тока ($ \vec{I}_\textit{макс} $).

\begin{align}
\begin{split}
\vec{I}_1&=\vec{I}_\textit{макс}\\
\vec{I}_2&=\sum_{n=1}^N \vec{I}_{n} - \vec{I}_\textit{макс}
\end{split}
\end{align}

\item
Выбор присоединения с максимальным током основана на сравнении мгновенных значений токов всех присоединений.
В момент начала аварийного режима на выходе логики фиксируется комплексное действующее значение основной гармоники тока того присоединения, у которого на этот момент зафиксирован максимальный ток.
Данное присоединение определяется как поврежденное и его номер сохраняется на всё время действия КЗ.

\item
БВКЗ анализирует угол между двумя сформированными векторами.
При внешнем КЗ угол между векторами превышает 90º.
Данный признак используется как критерий для формирования сигнала блокировки.

\medskip

\begin{figure}[H]
\centering
\begin{subfigure}{0.45\textwidth}
    \centering
    \includegraphics[width=\textwidth, trim=0 0 0 0, clip]{\fbpath/661.0631 ДЗШ/img/BVKZ_a.pdf}
    \caption{а) внутреннее КЗ}
    \label{pict:Action_DFZ_a}
\end{subfigure}
%\hfill
\begin{subfigure}{0.45\textwidth}
    \centering
    \includegraphics[width=\textwidth, trim=0 0 0 0, clip]{\fbpath/661.0631 ДЗШ/img/BVKZ_b.pdf}
    \caption{б) внешнее КЗ}
    \label{pict:Action_DFZ_b}
\end{subfigure}

\caption{Характеристика срабатывания БВКЗ}
\label{pict:BVKZ}
\end{figure}

\medskip

\item
Действие БВКЗ блокируется в режиме опробования шин, поскольку в данном режиме контролируется ток только одного присоединения, что исключает возможность корректного сравнения направлений векторов.

\end{enumerate}


%%%%%%%%%%%%%%%%%%%%%%%%%%%%%%%%%%%%%%%%%%%%%%%%%% Чувствительный токовый орган %%%%%%%%%%%%%%%%%%%%%%%%%%%%%%%%%%%%%%%%%%%%%%%%%%
\item Чувствительный токовый орган

\begin{enumerate}[label=\arabic{section}.\arabic{subsection}.\arabic{enumi}.\arabic*, labelsep=4pt, leftmargin=0em, itemindent=65pt, parsep=0pt]

\item
В устройстве предусмотрено два подхода для надежного отключения выключателей присоединений при работе ДЗШ, в том числе в режимах АПВ шин и опробования.
Такие режимы могут сопровождаться пониженным уровнем токов КЗ, в результате чего чувствительность дифференциальных органов может быть недостаточной.
Первый подход подразумевает применение ЧТО. Второй подход заключается в очувствлении уставок тормозной характеристики, которые используются в ДОТХ, БДО и МДО.

\item
ЧТО выполнен в виде реле тока, включенного на дифференциальные токи пускового и избирательных органов.
Ток срабатывания ЧТО для каждого дифференциального органа задается отдельной уставкой \textbf{«Iчто ПО»} (\textbf{«Iчто ИО1»}, \textbf{«Iчто ИО2»}).

\item
Очувствление представляет собой автоматическую смену величины начального тока срабатывания тормозной характеристики на более чувствительную уставку.
Для каждого дифференциального органа параметры тормозной характеристики в режиме очувствления задаются отдельными уставками.
Подробное описание алгоритма очувствления приведен в разделе «Логика очувствления».

\end{enumerate}


%%%%%%%%%%%%%%%%%%%%%%%%%%%%%%%%%%%%%%%%%%%%%%%%%% Блокировка по второй гармонике дифференциального тока %%%%%%%%%%%%%%%%%%%%%%%%%%%%%%%%%%%%%%%%%%%%%%%%%%
\item Блокировка по второй гармонике дифференциального тока

\begin{enumerate}[label=\arabic{section}.\arabic{subsection}.\arabic{enumi}.\arabic*, labelsep=4pt, leftmargin=0em, itemindent=65pt, parsep=0pt]

\item
Для предотвращения срабатывания ИО ДЗШ (за исключением ДТО) при БНТ, возникающих при включении присоединения с силовым трансформатором, используется блокировка по второй гармонике дифференциального тока.
Влияние второй гармоники, возникающей при БНТ, на работу ИО ДЗШ проявляется только в условиях опробования по схеме с «открытым плечом».

\item
ИО блокировки выполнен в виде реле тока, реагирующее на отношение дифференциального тока второй гармоники к дифференциальному току первой гармоники.
Величина срабатывания для каждого дифференциального органа задается отдельной уставкой \textbf{«I2г/I1г ПО»} (\textbf{«I2г/I1г ИО1»}, \textbf{«I2г/I1г ИО2»}).

\item
При срабатывании ИО ДЗШ сигнал на отключение сформируется по истечению соответствующей выдержки времени \textbf{«Т ДЗШ ПО»} (\textbf{«Т ДЗШ ИО1»}, \textbf{«Т ДЗШ ИО2»}).
Уставка времени срабатывания задается отдельно для каждого дифференциального органа.

\item
При нарушении фиксации присоединений сигналы отключения от ДЗШ и ЧТО формируются для обеих СШ.

\medskip

\begin{figure}[H]
\centering
\includegraphics[width=0.9\textwidth,height=0.9\textheight,keepaspectratio]{\fbpath/661.0631 ДЗШ/img/Gen_log.pdf}
\captionof{figure}{Функциональный блок ДЗШ}\label{pict:Gen_log}
\end{figure}

\medskip

\end{enumerate}


%%%%%%%%%%%%%%%%%%%%%%%%%%%%%%%%%%%%%%%%%%%%%%%%%% Контроль цепей тока %%%%%%%%%%%%%%%%%%%%%%%%%%%%%%%%%%%%%%%%%%%%%%%%%%
\item Контроль цепей тока

\begin{enumerate}[label=\arabic{section}.\arabic{subsection}.\arabic{enumi}.\arabic*, labelsep=4pt, leftmargin=0em, itemindent=65pt, parsep=0pt]

\item
В устройстве реализован контроль исправности токовых цепей для исключения ложного срабатывания ДЗШ вследствие нарушения баланса токов в дифференциальной цепи при обрыве, замыканий и т.д.

\medskip

\begin{figure}[H]
\centering
\includegraphics[width=0.9\textwidth,height=0.9\textheight,keepaspectratio]{\fbpath/661.0631 ДЗШ/img/KCT.pdf}
\captionof{figure}{Функциональный блок КЦТ}\label{pict:KCT}
\end{figure}

\medskip

\item
Ввод и вывод контроля исправности токовых цепей для каждого дифференциального органа осуществляется отдельным программным ключом \textbf{«КЦТ ПО»} (\textbf{«КЦТ ИО1»}, \textbf{«КЦТ ИО2»}).
В режиме нарушенной фиксации КЦТ избирательных органов выводится из работы.

\item
Для выявления неисправности токовых цепей предусмотрено три измерительных органа, реагирующих на дифференциальные токи ПО, ИО1 и ИО2.
Ток срабатывания КЦТ для каждого дифференциального органа задается отдельной уставкой \textbf{«Iкцт ПО»} (\textbf{«Iкцт ИО1»}, \textbf{«Iкцт ИО2»}).

\item
При обнаружении небаланса в токовых цепях измерительный орган с соответствующей выдержкой времени \textbf{«Ткцт ПО»} (\textbf{«Ткцт ИО1»}, \textbf{«Ткцт ИО2»}) действует на сигнализацию, а также на блокировку ДЗШ.
Блокировка ДЗШ от КЦТ осуществляется при введенном программном ключе \textbf{«Реж. КЦТ ПО»} (\textbf{«Реж. КЦТ ИО1»}, \textbf{«Реж. КЦТ ИО2»}).
Данные программные ключи приведены в перечне уставок функции ДЗШ.

\item
Сброс блокировки выполняется внешним сигналом \textit{«Возврат КЦТ»}.
Оперативный вывод блокировки ДЗШ при неисправностях в токовых цепях осуществляется внешним сигналом \textbf{«Вывод блок. ДЗШ от КЦТ»}.


\small
\begin{longtable}{|>{\centering\arraybackslash}m{5.3cm}|>{\centering\arraybackslash}m{3.3cm}|>{\centering\arraybackslash}m{4.2cm}|>{\centering\arraybackslash}m{1.8cm}|>{\centering\arraybackslash}m{1cm}|}
% Первая страница: полный заголовок
\caption{Перечень уставок функции ДЗШ\hfill\vspace{-0.5\baselineskip}}\label{desc:tbl2}\\
\hline
\rowcolor{gray!30}
Параметр & Условное обозначение на схеме & Значение/ Диапазон & Единица измерения & Шаг \\
\hline
\endfirsthead

% Промежуточные страницы: «Продолжение таблицы»
\caption*{\hspace{3pt}\emph{Продолжение таблицы \ref{desc:tbl2}}\hfill\ }\vspace{-0.5\baselineskip} \\
\hline
\rowcolor{gray!30}
Параметр & Условное обозначение на схеме & Значение/ Диапазон & Единица измерения & Шаг \\
\hline
\endhead

%===m>ДЗШ|-
%===m
\end{longtable}
\normalsize	

\end{enumerate}


%%%%%%%%%%%%%%%%%%%%%%%%%%%%%%%%%%%%%%%%%%%%%%%%%% Логика управления топологией %%%%%%%%%%%%%%%%%%%%%%%%%%%%%%%%%%%%%%%%%%%%%%%%%%
\item Логика управления топологией

\begin{enumerate}[label=\arabic{section}.\arabic{subsection}.\arabic{enumi}.\arabic*, labelsep=4pt, leftmargin=0em, itemindent=65pt, parsep=0pt]

\item
Логика управления топологией предназначена для назначения фиксации присоединений к СШ в соответствии со схемой первичных соединений.
С целью учета различных особенностей выполнения схем СШ в устройствах все присоединения имеют одинаковые возможности в задании конфигурации, т.е. каждое присоединение может быть зафиксировано за первой или второй СШ или назначено в качестве ШСВ посредством программного переключателя \textbf{«Фиксация Вn»}, где n – номер присоединения:
\begin{itemize} 
\item состояние \textbf{«Внеш.»} применяется для фиксации присоединения по состоянию двух входных сигналов: «Вn к 1С» и «Вn к 2С». В качестве внешних сигналов, например, могут выступать положения блок-контактов разъединителей на присоединении.
При пропадании срабатывания внешних сигналов логика сохраняет последнюю фиксацию присоединения. Сброс фиксации происходит при отсутствии состояния \textbf{«Внеш.»} (переход на другое состояние) или при срабатывании смежного внешнего сигнала.
\item состояние \textbf{«Вывод»} применяется для вывода присоединения из схемы ДЗШ. В этом состоянии ток присоединения не участвует в расчете дифференциальных и тормозных токов, используемых дифференциальными органами.
При выводе выключателя присоединения в ремонт (при поступлении внешнего сигнала \textit{«Ремонт Вn»}) ток присоединения также исключается из расчетов дифференциальных и тормозных токов.
\item состояние \textbf{«1С»} применяется для жесткой фиксации присоединения за первой СШ. В этом состоянии ток присоединения участвует в расчете дифференциальных и тормозных токов ИО1 и ПО.
\item состояние \textbf{«2С»} применяется для жесткой фиксации присоединения за второй СШ. В этом состоянии ток присоединения участвует в расчете дифференциальных и тормозных токов ИО2 и ПО.
\item состояние \textbf{«ШСВ1»} устанавливается для присоединения типа ШСВ в схеме ДЗШ с двумя ТТ ШСВ. В это состоянии фиксируется ТТ, установленный со стороны второй СШ. Присоединение учитывается за первой СШ, и его ток включается в расчет дифференциального и тормозного токов ИО1.
\item состояние \textbf{«ШСВ2»} устанавливается для присоединения типа ШСВ в схеме ДЗШ с двумя ТТ ШСВ. В это состоянии фиксируется ТТ, установленный со стороны первой СШ. Присоединение учитывается за второй СШ, и его ток включается в расчет дифференциального и тормозного токов ИО2.
\item состояние \textbf{«ШСВ12»} используется в схемах ДЗШ с одним ТТ ШСВ, установленного со стороны второй СШ. В этом состоянии присоединение учитывается за первой и второй СШ: ток включается в расчет дифференциальных и тормозных токов ИО1 со знаком «плюс», а в расчет токов ИО2 – со знаком «минус».
\item состояние \textbf{«ШСВ21»} используется в схемах ДЗШ с одним ТТ ШСВ, установленного со стороны первой СШ. В этом состоянии присоединение учитывается за первой и второй СШ: ток включается в расчет дифференциальных и тормозных токов ИО2 со знаком «плюс», а в расчет токов ИО1 – со знаком «минус».
\end{itemize}

\medskip

\begin{figure}[H]
\centering
\includegraphics[width=0.9\textwidth,height=0.9\textheight,keepaspectratio]{\fbpath/661.0631 ДЗШ/img/Fix_Q1.pdf}
\captionof{figure}{Функциональный блок органа фиксации В1}\label{pict:Fix_Q1}
\end{figure}

\medskip

\item
Режим нарушенной фиксации инициируется либо оперативно – посредством внешнего сигнала \textit{«Нарушение фикс.»}, либо автоматически – через сигналы, указывающие на нарушение фиксации конкретного присоединения.

\item
Сигнал нарушения фиксации отдельного присоединения может быть сформирован в следующих случаях:
\begin{itemize}
\item при отсутствии сигнала \textit{«Ремонт Вn»}, когда выключатель присоединения не находится в ремонте, фиксация осуществляется от внешних сигналов (положение \textbf{«Внеш.»}), и срабатывают оба сигнала фиксации присоединений: \textit{«Вn к 1С»} и \textit{«Вn к 2С»};
\item при фиксации присоединений по внешним сигналам (положение \textbf{«Внеш.»}), если в данном состоянии внешние сигналы фиксации исчезнут, но по присоединению продолжит протекать ток.
\end{itemize}

\item
Выходной сигнал \textit{«Нарушение фикс.»} удерживается в течение 200 мс после восстановление фиксации.
Данная выдержка необходима для исключения влияния переходных состояний, возникающих при переключении разъединителей или кратковременном пропадании входного сигнала \textit{«Вn к 1С»} (\textit{«Вn к 2С»}).

\medskip

\begin{figure}[H]
\centering
\includegraphics[width=0.9\textwidth,height=0.9\textheight,keepaspectratio]{\fbpath/661.0631 ДЗШ/img/LUT.pdf}
\captionof{figure}{Функциональный блок логики нарушение фиксации}\label{pict:LUT}
\end{figure}

\medskip

\item
Уставки функции фиксации присоединений представлены в таблице \ref{desc:tbl3}.

\small
\begin{longtable}{|>{\centering\arraybackslash}m{5.3cm}|>{\centering\arraybackslash}m{3.3cm}|>{\centering\arraybackslash}m{4.2cm}|>{\centering\arraybackslash}m{1.8cm}|>{\centering\arraybackslash}m{1cm}|}
% Первая страница: полный заголовок
\caption{Перечень уставок функции фиксации присоединений\hfill\vspace{-0.5\baselineskip}}\label{desc:tbl3}\\
\hline
\rowcolor{gray!30}
Параметр & Условное обозначение на схеме & Значение/ Диапазон & Единица измерения & Шаг \\
\hline
\endfirsthead

% Промежуточные страницы: «Продолжение таблицы»
\caption*{\hspace{3pt}\emph{Продолжение таблицы \ref{desc_dzl:tbl3}}\hfill\ }\vspace{-0.5\baselineskip} \\
\hline
\rowcolor{gray!30}
Параметр & Условное обозначение на схеме & Значение/ Диапазон & Единица измерения & Шаг \\
\hline
\endhead

%===m>ДЗШ|-
%===m
\end{longtable}
\normalsize	


\end{enumerate}

%%%%%%%%%%%%%%%%%%%%%%%%%%%%%%%%%%%%%%%%%%%%%%%%%% Контроль цепей напряжения %%%%%%%%%%%%%%%%%%%%%%%%%%%%%%%%%%%%%%%%%%%%%%%%%%
\item Контроль цепей напряжения

\begin{enumerate}[label=\arabic{section}.\arabic{subsection}.\arabic{enumi}.\arabic*, labelsep=4pt, leftmargin=0em, itemindent=65pt, parsep=0pt]

\item
Логика контроля цепей напряжения предназначена:
\begin{itemize}
\item для диагностики неисправности во вторичных цепях ТН;
\item для контроля отсутствия напряжения на СШ для ввода очувствления ДЗШ при опробовании;
\item для контроля наличия напряжения на СШ для сброса автоматического очувствления ДЗШ при успешном АПВ шин;
\item для запрета АПВ при наличии напряжения на СШ (при неполнофазном или полнофазном отказе выключателя питающего присоединения).
\end{itemize}

\medskip

\begin{figure}[H]
\centering
\includegraphics[width=0.9\textwidth,height=0.9\textheight,keepaspectratio]{\fbpath/661.0631 ДЗШ/img/KCN.pdf}
\captionof{figure}{Функциональный блок КЦН 1С}\label{pic:KCN}
\end{figure}

\medskip

\item
Ввод и вывод в работу контроля исправности цепей напряжения для каждого ТН на СШ осуществляется либо посредством программного ключа \textbf{«КЦН 1С»} (\textbf{«КЦН 2С»}), либо с использованием внешнего сигнала \textit{«Вывод КЦН 1С»} (\textit{«Вывод КЦН 2С»}).

\item
При снижении значения хотя бы одного линейного напряжения ниже уставки \textbf{«ИО отсутствия U 1C»} (\textbf{«ИО отсутствия U 2C»}), либо при срабатыванием ИО наличия напряжения ОП (\textbf{«ИО наличия ОП U 1C»}, \textbf{«ИО наличия ОП U 2C»}) через выдержку времени \textbf{«Ткцн 1С»} (\textbf{«Ткцн 2С»}) формируется сигнал \textit{«Неиспр. ЦН 1С»} (\textit{«Неиспр. ЦН 2С»}).
Данный сигнал свидетельствует о наличии неисправности цепей напряжения ТН на СШ.

\item
Сигнал \textit{«Отсутствие U 1С»} (\textit{«Отсутствие U 2С»}) формируется, если значения всех линейных напряжений опускается ниже уставки \textbf{«ИО отсутствия U 1C»} (\textbf{«ИО отсутствия U 2C»}) и одновременно отсутствует срабатывание ИО наличия напряжения ОП.
Данный сигнал свидетельствует об отсутствии напряжения и переходного режима на СШ.

\item
Сигнал \textit{«Наличие U при откл. 1С»} (\textit{«Наличие U при откл. 2С»}) формируется, если значение хотя бы одного линейного напряжения превысят уставку \textbf{«ИО наличия U 1C»} (\textbf{«ИО наличия U 2C»}), либо при срабатывании ИО наличия напряжения ОП (\textbf{«ИО наличия ОП U 1C»}, \textbf{«ИО наличия ОП U 2C»}).
Данный сигнал свидетельствует о наличии напряжения на отключенной СШ и используется в логике запрета АПВ для факта наличия неполнофазного (полнофазного) отказа выключателя.

\item
Сигнал \textit{«Наличие U 1С»} (\textit{«Наличие U 2С»}) формируется, если значения всех линейных напряжений превысят уставку \textbf{«ИО наличия U 1C»} (\textbf{«ИО наличия U 2C»}) и одновременно отсутствует срабатывание ИО наличия напряжения ОП.
Данный сигнал предназначен для сброса автоматического очувствления после успешного АПВ первого присоединения и свидетельствует о нормальном режиме работы и отсутствии переходного режима на СШ.

\item
Уставки функции КЦН представлены в таблице \ref{desc:tbl4}.

\small
\begin{longtable}{|>{\centering\arraybackslash}m{5.3cm}|>{\centering\arraybackslash}m{3.3cm}|>{\centering\arraybackslash}m{4.2cm}|>{\centering\arraybackslash}m{1.8cm}|>{\centering\arraybackslash}m{1cm}|}
% Первая страница: полный заголовок
\caption{Перечень уставок функции КЦН\hfill\vspace{-0.5\baselineskip}}\label{desc:tbl4}\\
\hline
\rowcolor{gray!30}
Параметр & Условное обозначение на схеме & Значение/ Диапазон & Единица измерения & Шаг \\
\hline
\endfirsthead

% Промежуточные страницы: «Продолжение таблицы»
\caption*{\hspace{3pt}\emph{Продолжение таблицы \ref{desc:tbl4}}\hfill\ }\vspace{-0.5\baselineskip} \\
\hline
\rowcolor{gray!30}
Параметр & Условное обозначение на схеме & Значение/ Диапазон & Единица измерения & Шаг \\
\hline
\endhead

%===m>ДЗШ|-
%===m
\end{longtable}
\normalsize	

\end{enumerate}


%%%%%%%%%%%%%%%%%%%%%%%%%%%%%%%%%%%%%%%%%%%%%%%%%% Логика опробования %%%%%%%%%%%%%%%%%%%%%%%%%%%%%%%%%%%%%%%%%%%%%%%%%%
\item Логика опробования

\begin{enumerate}[label=\arabic{section}.\arabic{subsection}.\arabic{enumi}.\arabic*, labelsep=4pt, leftmargin=0em, itemindent=65pt, parsep=0pt]

\item
Логика опробования предназначена для работы ДЗШ при оперативном опробовании выключателей присоединений по схеме с «открытым» или «закрытым» плечом.
При неуспешном опробовании отключается только опробуемый выключатель, остальные выключатели остаются в работе.

\medskip

\begin{figure}[H]
\centering
\includegraphics[width=0.9\textwidth,height=0.9\textheight,keepaspectratio]{\fbpath/661.0631 ДЗШ/img/Log_oprob.pdf}
\captionof{figure}{Функциональный блок логики опробования}\label{pic:Log_oprob}
\end{figure}

\medskip

\item
Переход в режим опробования осуществляется при наличии разрешающего сигнала \textit{«Ввод опробования»}.
Запуск опробования выполняется при поступлении команды РКВ от выключателя любого присоединения, зафиксированного за определенной СШ (при наличии сигнала \textit{«Фиксация Вn 1С»} или \textit{«Фиксация Вn 2С»}, где n – номер присоединения).
Продолжительность опробования определяется уставкой выдержки времени \textbf{«Т опроб»}.

\item
С помощью программного переключателя \textbf{«Опроб. Вn»} задается условие для отключения опробуемого выключателя по срабатыванию одного из следующих органов: ДЗШ ПО, ЧТО ПО или индивидуального ИО по току опробуемого присоединения с блокировкой по второй гармонике.

\item
Ток срабатывания индивидуального ИО задается с помощью уставки \textbf{«I опроб. Вn»}.
Блокировка по второй гармонике необходима для предотвращения ложных срабатываний ИО при БНТ трансформатора.
Для ввода блокировки предусмотрен программный ключ \textbf{«БНТ Вn»}. Уровень срабатывания блокировки задается с помощью уставки \textbf{«I2г/ I1г Вn»}.
Срабатывание ИО осуществляется по завершении выдержки времени \textbf{«Тср. опроб. Bn»}.

\item
Для обеспечения отключения только опробуемого выключателя, без отключения всей СШ, предусмотрены сигналы блокировки ДЗШ: \textit{«Блок. ДЗШ 1С при опроб.»} (\textit{«Блок. ДЗШ 2С при опроб.»} и \textit{«Блок. ДЗШ при опроб.»}).
Время действия блокировки соответствует продолжительности опробования.

\item
При нарушении фиксации присоединений сигналы блокировки ДЗШ формируются для обеих СШ.

\small
\begin{longtable}{|>{\centering\arraybackslash}m{5.3cm}|>{\centering\arraybackslash}m{3.3cm}|>{\centering\arraybackslash}m{4.2cm}|>{\centering\arraybackslash}m{1.8cm}|>{\centering\arraybackslash}m{1cm}|}
% Первая страница: полный заголовок
\caption{Перечень уставок функции опробования\hfill\vspace{-0.5\baselineskip}}\label{desc:tbl5}\\
\hline
\rowcolor{gray!30}
Параметр & Условное обозначение на схеме & Значение/ Диапазон & Единица измерения & Шаг \\
\hline
\endfirsthead

% Промежуточные страницы: «Продолжение таблицы»
\caption*{\hspace{3pt}\emph{Продолжение таблицы \ref{desc:tbl5}}\hfill\ }\vspace{-0.5\baselineskip} \\
\hline
\rowcolor{gray!30}
Параметр & Условное обозначение на схеме & Значение/ Диапазон & Единица измерения & Шаг \\
\hline
\endhead

%===m>ДЗШ|-
%===m
\end{longtable}
\normalsize	

\end{enumerate}


%%%%%%%%%%%%%%%%%%%%%%%%%%%%%%%%%%%%%%%%%%%%%%%%%% Логика очувствления %%%%%%%%%%%%%%%%%%%%%%%%%%%%%%%%%%%%%%%%%%%%%%%%%%
\item Логика очувствления

\begin{enumerate}[label=\arabic{section}.\arabic{subsection}.\arabic{enumi}.\arabic*, labelsep=4pt, leftmargin=0em, itemindent=65pt, parsep=0pt]

\item
В устройстве предусмотрено два режима очувствления ДЗШ: оперативный и автоматический.
Оперативный режим вводится наличием внешнего сигнала \textit{«Оперативный ввод очувствления»} и применяется для обеспечения требуемой чувствительности ДЗШ при выводе в ремонт мощных присоединений.
Автоматический режим вводится наличием внешнего сигнала \textit{«Нормальный режим очувствления»} или от программного ключа \textbf{«Очувст.»} и предназначен для повышения надежности отключения выключателей присоединений при работе ДЗШ, в том числе при АПВ шин и при опробовании.

\item
Режим очувствления, вводимый срабатыванием ДЗШ, сохраняется на всё время отключения СШ.
При наличии сигнала \textit{«Фикс. откл 1С»} (\textit{«Фикс. откл. 2С»}) режим очувствления удерживается до момента возможного повторного включения СШ при АПВ.
Сброс очувствления до повторного включения осуществляется при наличии сигнала \textit{«Работа АПВ 1С»} (\textit{«Работа АПВ 2С»}).

\item
Переход в режим очувствления при опробовании осуществляется при отсутствии напряжения на СШ.
Очувствление в данном случае сохраняется на всё время опробования.

\item
При АПВ шин режим очувствления вводится на время включения выключателя и последующего отключения в случае неуспешного АПВ.
Для исключения излишнего очувствления при постепенном включении всех выключателей предусмотрена возможность вывода режима очувствления при успешном АПВ по истечению выдержки времени, задаваемой уставкой \textbf{«Точ. АПВ 1С»} (\textbf{«Точ. АПВ 2С»}).
Вывод режима осуществляется программным ключом \textbf{«Вывод очув. при успеш. АПВ 1С»} (\textbf{«Вывод очув. при успеш. АПВ 2С»}).

\medskip

\begin{figure}[H]
\centering
\includegraphics[width=0.9\textwidth,height=0.9\textheight,keepaspectratio]{\fbpath/661.0631 ДЗШ/img/Log_ochuv.pdf}
\captionof{figure}{Функциональный блок логики очувствления}\label{pic:Log_ochuv}
\end{figure}

\medskip

\small
\begin{longtable}{|>{\centering\arraybackslash}m{5.3cm}|>{\centering\arraybackslash}m{3.3cm}|>{\centering\arraybackslash}m{4.2cm}|>{\centering\arraybackslash}m{1.8cm}|>{\centering\arraybackslash}m{1cm}|}
% Первая страница: полный заголовок
\caption{Перечень уставок функции логики очувствления\hfill\vspace{-0.5\baselineskip}}\label{desc:tbl6}\\
\hline
\rowcolor{gray!30}
Параметр & Условное обозначение на схеме & Значение/ Диапазон & Единица измерения & Шаг \\
\hline
\endfirsthead

% Промежуточные страницы: «Продолжение таблицы»
\caption*{\hspace{3pt}\emph{Продолжение таблицы \ref{desc:tbl6}}\hfill\ }\vspace{-0.5\baselineskip} \\
\hline
\rowcolor{gray!30}
Параметр & Условное обозначение на схеме & Значение/ Диапазон & Единица измерения & Шаг \\
\hline
\endhead

%===m>ДЗШ|-
%===m
\end{longtable}
\normalsize	

\end{enumerate}


%%%%%%%%%%%%%%%%%%%%%%%%%%%%%%%%%%%%%%%%%%%%%%%%%% Требования к цифровому каналу связи ДЗЛ %%%%%%%%%%%%%%%%%%%%%%%%%%%%%%%%%%%%%%%%%%%%%%%%%%
\item Требования к цифровому каналу связи ДЗЛ

\begin{enumerate}[label=\arabic{section}.\arabic{subsection}.\arabic{enumi}.\arabic*, labelsep=4pt, leftmargin=0em, itemindent=65pt, parsep=0pt]

\item
Для корректной работы ДЗЛ канал связи (с учетом работы оборудования ЦСПИ) должен соответствовать следующим требованиям:
\begin{itemize} 
\item максимальная допустимая задержка передачи данных в одном направлении не более 30 мс;
\item максимальная допустимая асимметрия задержки передачи данных (разность времени передачи данных в прямом и обратном направлениях) не более 0,5 мс;
\item одновременное в обоих направлениях (на прием и на передачу) переключение КС на резервный маршрут при потере основного маршрута или обеспечение статического маршрута без автоматического переключения при передаче данных через мультиплексированные КС.
\end{itemize}

\end{enumerate}


%%%%%%%%%%%%%%%%%%%%%%%%%%%%%%%%%%%%%%%%%%%%%%%%%% Параметры настройки цифрового канала связи ДЗЛ %%%%%%%%%%%%%%%%%%%%%%%%%%%%%%%%%%%%%%%%%%%%%%%%%%
\item Параметры настройки цифрового канала связи ДЗЛ

\begin{enumerate}[label=\arabic{section}.\arabic{subsection}.\arabic{enumi}.\arabic*, labelsep=4pt, leftmargin=0em, itemindent=65pt, parsep=0pt]

\item
Параметр \textbf{«Синх. КС»} задает источник синхронизации передаваемых данных в КС и должен быть определен исходя из его типа (выделенный или мультиплексированный).
Для синхронизации данных может быть использован синхросигнал от внутреннего тактового генератора (Внутренняя) либо синхронизирующий сигнал, выделяемый из принимаемых данных локального п/к ДЗЛ (Внешняя).
В случае применения мультиплексированных КС параметр \textbf{«Синх. КС»} должен устанавливаться в значение \textit{Внешняя} независимо от способа подключения к мультиплексору.
Обязательным требованием к мультиплексированному КС является синхронность по каналу передачи и каналу приема, что должно обеспечиваться соответствующими настройками оборудования связи.
В случае применения выделенных КС отсутствует необходимость использовать внешнюю синхронизацию по причине отсутствия дополнительных устройств в канале связи, с которыми необходимо синхронизировать данные.
Поэтому параметр \textbf{«Синх. КС»} должен быть установлен в значение \textit{Внутренняя}.
При использовании кодирования, предусмотренного стандартом С37.94 на выделенном канале связи возможно использование как внутренней, так и внешней синхронизации.
Таким образом один из п/к ДЗЛ может быть источником сигнала синхронизации, а другой п/к ДЗЛ может быть приемником сигнала синхронизации.

\item
Параметр \textbf{«Скор»} на выделенных КС устанавливает логическую скорость при использовании внутренней синхронизации.
При использовании внешней синхронизации параметр игнорируется.
При использовании мультиплексированного КС параметр \textbf{«Скор»} игнорируется и может быть выставлен любым, так как скорость передачи данных выставляется непосредственно в мультиплексоре установкой числа N, которое определяет количество 64-килобитных каналов, используемых для формирования общей скорости передачи.
По умолчанию параметр \textbf{«Скор»} равен 64 кБит/с. При наличии двух плеч на присоединение параметр \textbf{«Скор»} должен быть равен не менее 128 кБит/с.

\item
Параметр \textbf{«Комп. асим.»} необходим для компенсации постоянной составляющей асимметрии КС.
На выделенных КС параметр \textbf{«Комп. асим.»} устанавливается равным 0, так как длина ВОЛС в прямом и обратном направлении, как правило, одинакова.
Поэтому выделенные КС принципиально симметричны.
Для мультиплексированных КС необходимо задать режим п/к ДЗЛ уставкой \textbf{«Реж. терм.»} При этом необходимо правильно устанавливать режим каждого п/к ДЗЛ.
Например, в случае установки в обоих п/к ДЗЛ режима «Ведущий» каждый п/к ДЗЛ использует собственное установленное время для компенсации асимметрии.
В случае установки одинакового режима в обоих п/к ДЗЛ компенсация асимметрии будет работать неправильно.
П/к ДЗЛ с режимом «Ведущий» использует заданное значение параметра \textbf{«Комп. асим.»} и передает данное значение в удаленный п/к ДЗЛ.
П/к ДЗЛ с режимом «Ведомый» игнорирует значение параметра \textbf{«Комп. асим.»} и использует принятое значение от удаленного п/к ДЗЛ с противоположным знаком.
Рекомендации по определению параметра компенсации асимметрии КС приведены {\color{red}в приложении А}.

\item
Параметры \textbf{«ID Tx»} и \textbf{«ID Rx»} предназначены для исключения ошибок при подключении каналов связи вследствие неправильной коммутации потоков данных.
Использование уникальных идентификаторов для каждого КС (для приема и передачи данных) позволяет выявить подобные ошибки и исключить неправильную работу функции ДЗЛ.
П/к ДЗЛ проверяет соответствие установленного параметра \textbf{«ID Rx»} и принимаемого идентификатора в данных от удаленного п/к ДЗЛ.
В случае их несоответствия формируется сигнал неисправности канала связи.

\end{enumerate}


%%%%%%%%%%%%%%%%%%%%%%%%%%%%%%%%%%%%%%%%%%%%%%%%%% Диагностика состояния цифрового канала связи %%%%%%%%%%%%%%%%%%%%%%%%%%%%%%%%%%%%%%%%%%%%%%%%%%
\item Диагностика состояния цифрового канала связи

\begin{enumerate}[label=\arabic{section}.\arabic{subsection}.\arabic{enumi}.\arabic*, labelsep=4pt, leftmargin=0em, itemindent=65pt, parsep=0pt]

\item
Каждый п/к ДЗЛ осуществляет автоматический непрерывный контроль состояния каналов связи.
Для каждого канала связи (ЦКС1 и ЦКС2) осуществляется регистрация в энергонезависимой памяти терминала и отображение на ИЧМ терминала диагностических сигналов для контроля исправности, а также проверки правильности задания параметров настройки. 

\item
В таблице ниже приведен перечень диагностических сигналов ЦКС1, для ЦКС2 перечень аналогичный.
\small
\begin{longtable}{|>{\centering\arraybackslash}m{4cm}|>{\centering\arraybackslash}m{6.4cm}|>{\centering\arraybackslash}m{2.8cm}|>{\centering\arraybackslash}m{2.8cm}|}
% Первая страница: полный заголовок
\caption{Перечень диагностических сигналов ЦКС1\hfill\vspace{-0.5\baselineskip}}\label{desc_dzl:tbl10}\\
\hline
\rowcolor{gray!30}
Обозначение & Описание & Знач. при исправном состоянии & Знач. при неисправ. состоянии \\
\hline
\endfirsthead

% Промежуточные страницы: «Продолжение таблицы»
\caption*{\hspace{3pt}\emph{Продолжение таблицы \ref{desc_dzl:tbl10}}\hfill\ }\vspace{-0.5\baselineskip} \\
\hline
\rowcolor{gray!30}
Обозначение & Описание & Знач. при исправном состоянии & Знач. при неисправ. состоянии \\
\hline
\endhead

\centering\arraybackslash ЦКС1: Приним. ID  & \centering\arraybackslash Принимаемый идентификатор устройства ЦКС1 & \centering\arraybackslash =ID Rx & \centering\arraybackslash ≠ ID Rx \\
\hline
\centering\arraybackslash ЦКС1: t\_ПЕР  & \centering\arraybackslash Текущее значение времени задержки передачи данных по ЦКС1, мкс & \centering\arraybackslash 0-30000 & \centering\arraybackslash >30000 \\
\hline
\centering\arraybackslash ЦКС1: t\_АС  & \centering\arraybackslash Измеренное значение времени асимметрии задержки передачи данных по ЦКС1 (при наличии внешней синхронизации), мкс & \centering\arraybackslash 0 & \centering\arraybackslash >|±500| \\
\hline
\centering\arraybackslash ЦКС1: Сбой синхр.  & \centering\arraybackslash Неисправность внешнего источника синхронизации  & \centering\arraybackslash 0 & \centering\arraybackslash 1 \\
\hline
\centering\arraybackslash ЦКС1: Уд. сбой синхр.  & \centering\arraybackslash Неисправность внешнего источника синхронизации на удаленном устройстве & \centering\arraybackslash 0 & \centering\arraybackslash 1 \\
\hline
\centering\arraybackslash ЦКС1: Отсут. приема  & \centering\arraybackslash Отсутствие приема данных по ЦКС & \centering\arraybackslash 0 & \centering\arraybackslash 1 \\
\hline
\centering\arraybackslash ЦКС1: Уд. отсут. приема  & \centering\arraybackslash Отсутствие приема данных по ЦКС1 на удаленном устройстве & \centering\arraybackslash 0 & \centering\arraybackslash 1 \\
\hline
\centering\arraybackslash ЦКС1: Неверный ID  & \centering\arraybackslash Прием данных с неверным адресом источника по ЦКС1 & \centering\arraybackslash 0 & \centering\arraybackslash 1 \\
\hline
\centering\arraybackslash ЦКС1: Сбой данных  & \centering\arraybackslash Неверная контрольная сумма принятых данных по ЦКС1 & \centering\arraybackslash 0 & \centering\arraybackslash 1 \\
\hline
\centering\arraybackslash ЦКС1: Сбой передачи  & \centering\arraybackslash Недопустимая задержка передачи данных по ЦКС1 & \centering\arraybackslash 0 & \centering\arraybackslash 1 \\
\hline
\centering\arraybackslash ЦКС1: Сч. ошибок  & \centering\arraybackslash Текущее количество ошибок в ЦКС1 & \centering\arraybackslash 0--10000, не изменяется & \centering\arraybackslash Увелич. на 1 \\
\hline
\centering\arraybackslash ЦКС1: Неисправность КС  & \centering\arraybackslash Неисправное состояние ЦКС1 & \centering\arraybackslash 0 & \centering\arraybackslash 1 \\
\hline

\end{longtable}
\normalsize

\item
Параметр \textit{Приним. ID} показывает идентификатор устройства КС, получаемый в данных от удаленного п/к ДЗЛ.
При исправном состоянии указанный параметр соответствует заданному значению \textbf{«ID Rx»}.

\item
Параметр \textit{t\_ПЕР} показывает измеряемое время задержки передачи данных в одном направлении, определяемое скоростью передачи данных, длиной канала связи, а также для мультиплексированных КС количеством мультиплексоров и способом переприема данных в них.
При исправном состоянии указанный параметр не превышает 30 мс. 

\item
Параметр \textit{t\_АС} показывает измеряемое время асимметрии задержки передачи данных в прямом и обратном направлении.
При отсутствии сигнала 1PPS хотя бы в одном из п/к ДЗЛ указанный параметр должен устанавливаться равным нулю, что не является признаком симметричности КС.
Измерение может осуществляться только при наличии сигналов 1PPS от устройств синхронизации GPS/ГЛОНАСС в обоих п/к ДЗЛ.
При исправном состоянии указанный параметр не превышает 0,5 мс. Значения отображаются с учетом выставленной уставки компенсации асимметрии \textbf{«Комп. асим.»}.

\item
Сигнал неисправности \textit{«Отсут. приема»} формируется при отсутствии данных на приемнике порта локального п/к ДЗЛ от удаленного устройства (п/к ДЗЛ или оборудования ЦСПИ). 

\item
Сигнал неисправности \textit{«Уд. отсут. приема»} формируется при отсутствии данных на приемнике порта удаленного устройства (п/к ДЗЛ или оборудования ЦСПИ) от локального п/к ДЗЛ.

\item
Сигнал неисправности \textit{«Неверный ID»} формируется при несоответствии принимаемого идентификатора КС удаленного п/к ДЗЛ заданному значению \textbf{«ID Rx»}.

\item
Сигнал неисправности \textit{«Сбой данных»} формируется при неверной контрольной сумме принятых данных, что означает превышение вероятности битовых ошибок в канале связи над требуемым значением для штатной работы п/к ДЗЛ (для канала связи через ЦСПИ - $ 10^{-3}) $.

\item
Сигнал неисправности \textit{«Сбой передачи»} формируется при превышении времени задержки передачи данных максимально допустимого значения, равного 30 мс.

\item
Сигнал \textit{«Сч. Ошибок»} не является признаком неисправности и показывает общее количество ошибок в КС, возникших до текущего момента времени.
При фиксации одной из ошибок в КС значение счетчика ошибок увеличивается на единицу.
Значение счетчика ошибок не сбрасывается при перерыве питания терминала.

\item
Общий сигнал неисправности КС формируется при фиксации одной из неисправностей в КС (\textit{«Отсут. приема», «Уд. отсут. Приема», «Неверный ID», «Сбой данных», «Сбой передачи», «Сбой синхр.», «Уд. сбой синхр.»}) в течение фиксированной выдержки времени 50 мс или при наличии соответствующего входного внешнего сигнала неисправности КС.

\end{enumerate}


\end{enumerate}
